\section{External validity in pre-committed novel settings}
\label{sec:ext-valid}

We now turn to making guarantees about inference in novel settings.
This is made possible when we have a pre-committed family of settings from which we can randomly sample.
In particular, this framework allows us to compare the relative accuracy of two models in predicting human responses even for unsampled settings---where we have no previous data---within the same general domain.

The setup is related to that of \cite{Hotz2005PredictingEfficacy}, \cite{Allcott2015Site}, and \cite{Dehejia2019External}, where treatment effects from various ``sites'' are used to evaluate the external validity of a given intervention at the population level.
However, they assume a common underlying intervention---analogous to a single setting in our framework---across all sites.
When there are heterogeneous interventions, only special instances with strong additional assumptions allow for appropriate inference.
The following, which is very similar to the theoretical framework in \cite{andrews2025transfer}, requires no such assumptions.

Let $X=\{x_{1},\dots,x_{M}\}$ denote the pool of candidate settings for which we wish to make predictions.
$P(y|x)$ denotes the true human response distribution for $y \in Y$---the set of allowable responses.
We define a predictive distribution for some flexible model $\theta$---an LLM, an economic model, etc.---as $\hat P_{\theta}(y|x)$.
For a given setting $x$, the expected log-likelihood that the human distribution could have been produced by $\theta$ is
\[
\ell(x;\theta) = \mathbb{E}_{y\sim P(\cdot\mid x)}\bigl[\log\hat P_{\theta}(y|x)\bigr].
\]
Then the comparative predictive power of two models $\theta'$ and $\theta''$ can be measured via
\begin{equation} \label{eq:ratio}
  \Lambda(x) = \ell(x;\theta') - \ell(x;\theta'').
\end{equation}
A positive $\Lambda(x)$ means that $\theta'$ assigns more probability mass to the human responses than $\theta''$ for $x$.
Averaging over all settings in $X$ yields the population estimand
\begin{equation} \label{eq:avg-ratio}
\bar{\Lambda} = \mathbb{E}_{x\sim\pi}\bigl[\Lambda(x)\bigr]
\end{equation}
where $\pi$ is some distribution over $X$.
A positive $\bar{\Lambda}$ is interpreted as evidence that $\theta'$ is, on average, more predictive of human behavior than $\theta''$ across the entire population of $X$.

Suppose we observe a sample $S=\{s_{1},\dots,s_{n}\}\subset X$ and, for each
$s \in S$, independent human responses $y_s=(y_{s,1},\dots,y_{s,m_s})$.
Identification to estimate equations~\ref{eq:ratio} and \ref{eq:avg-ratio} from these samples requires the following assumptions.

\begin{assumption}  \label{ass:rand-settings}
(Unconfounded settings).
The observed settings $S=\{s_{1},\dots,s_{n}\}$ are randomly sampled from distribution $\pi$ with full support over $X$ such that $\pi(x)>0$ for all $x \in X$.
\end{assumption}

\begin{assumption}\label{ass:rand-humans}
(Random assignment and within-setting independence).
Humans are randomly assigned to settings in $S$.
Human responses within a setting are independent draws from $P(y \mid s)$.
\end{assumption}

\begin{assumption} \label{ass:response-support}
(Positivity and finite second moment).
Whenever $P(y\mid x)>0$, then $\hat P_{\theta'}(y\mid x)>0$ and $\hat P_{\theta''}(y\mid x)>0$. 
Moreover, $\mathbb{E}_{x\sim\pi}\bigl\{\mathbb{E}_{y\sim P(\cdot\mid x)}\bigl[(\log\hat P_{\theta'}(y\mid x) - \log\hat P_{\theta''}(y\mid x))^{2}\bigr]\bigr\} < \infty$.
\end{assumption}

The first two assumptions are basically identical to the assumptions of \emph{unconfounded location} and \emph{random assignment} in \citeauthor{Hotz2005PredictingEfficacy}.
They are also very similar in spirit to Assumption 1 in \citeauthor{andrews2025transfer}.
The first part of Assumption~\ref{ass:response-support} is similar to the covariate overlap assumption in causal inference; without it, some observed responses would have $\log 0$.
The second portion of Assumption~\ref{ass:response-support} is a standard finite second-moment condition.

For every setting $s \in S$, define the sample analogue to equation~\ref{eq:ratio} as:
\begin{equation} \label{eq:sample-ratio}
\hat\Lambda_{s}
\;=\;
\frac1{m_s}
\sum_{j=1}^{m_s}
\Bigl[\log\hat P_{\theta'}(y_{s,j}\mid s)
-\log\hat P_{\theta''}(y_{s,j}\mid s)\Bigr].
\end{equation}
Aggregate across settings to produce the sample analogue to equation~\ref{eq:avg-ratio}:
\begin{equation} \label{eq:sample-avg-ratio}
\bar\Lambda_{S}
\;=\;
\frac1n\sum_{s \in S}\hat\Lambda_{s}.
\end{equation}

\begin{proposition}[Unbiasedness and asymptotic normality]
\label{prop:sample-valid}
Suppose Assumptions~\ref{ass:rand-settings}--\ref{ass:response-support} hold.
Then
\[
\mathbb{E}[\bar{\Lambda}_{S}] = \bar{\Lambda}
\quad \text{and} \quad
\sqrt{n}\,(\bar{\Lambda}_{S}-\bar{\Lambda})
\xrightarrow{d} \mathcal{N}\bigl(0,\sigma^{2}\bigr),
\]
where $
\sigma^{2}
= \operatorname{Var}_{x\sim\pi}\!\bigl[\Lambda(x)\bigr]
\;+\;
\mathbb{E}_{x\sim\pi}\!\Bigl[\tfrac{1}{m_x}\,V_x\Bigr]$ and 
$V_x
= \operatorname{Var}_{y \sim P(\cdot \mid x)}
\bigl[\log\hat P_{\theta'}(y \mid x)
-\log\hat P_{\theta''}( y\mid x)\bigr].$\footnote{\emph{Proof.} (Unbiasedness).
For any setting $s$, Assumption~\ref{ass:rand-humans} implies $\mathbb{E}\bigl[\hat\Lambda_s \mid s\bigr]=\Lambda(s).$
Hence $\mathbb{E}\bigl[\bar\Lambda_S \mid S\bigr] = \frac{1}{n}\sum_{s \in S}\Lambda(s).$
Taking expectation over the i.i.d.\ draw of the settings (Assumption~\ref{ass:rand-settings}) yields $\mathbb{E}[\bar\Lambda_S]=\bar\Lambda.$
(Asymptotic normality).
The random variables $\{\hat\Lambda_s\}_{s \in S}$ are i.i.d. across settings with finite variance $
\operatorname{Var}(\hat\Lambda_s)
= \operatorname{Var}_{x\sim\pi}[\Lambda(x)]
+ \mathbb{E}_{x\sim\pi}
\bigl[{\textstyle\frac{1}{m_x}}\,V_x\bigr]
\;<\;\infty,$
where the decomposition follows from the law of total variance and the independence of human draws within each setting.
Because the moment condition in Assumption~\ref{ass:response-support} guarantees $V_x<\infty$, the central-limit theorem gives $\sqrt{n}\bigl(\bar\Lambda_S-\bar\Lambda\bigr)
\;\xrightarrow{d}\;
\mathcal N\bigl(0,\sigma^{2}\bigr).$
(Regularity).
By Assumption~\ref{ass:response-support},
$\log\hat P_{\theta'}(y\mid x)$ and $\log\hat P_{\theta''}(y\mid x)$ are finite, so all moments used above exist.}
\end{proposition}

Notably, this does not rely on a large sample size of humans for any particular setting.
The asymptotic variance in Proposition~\ref{prop:sample-valid} decomposes into two conceptually distinct parts.
The first term, $\operatorname{Var}_{x\sim\pi}\bigl[\Lambda(x)\bigr]$, reflects heterogeneity in model performance across settings.
The second term, $\mathbb{E}_{x\sim\pi}\bigl[\frac{1}{m_x}V_x\bigr]$, is sampling noise that arises because we estimate $\Lambda(x)$ with a finite number $m_x$ of human draws. 
As long as $m_x\ge 1$, this component is finite. 
Consequently, once a few human observations are obtained per setting, the precision of $\bar\Lambda_S$ is governed primarily by $n$ 
because each setting contributes only a single observation $\hat\Lambda_s$.
This also means there is no within-cluster correlation left to adjust for, so the standard sample variance estimator ($\hat{\sigma}^{2}\;=\;\frac{1}{n-1}\sum_{s\in S}\!\bigl(\hat{\Lambda}_{s}-\bar{\Lambda}_{S}\bigr)^{2}$) across settings already yields valid standard errors without needing to adjust for clustering.
As such, standard $z$-tests or Wald confidence intervals follow immediately when estimating if $\bar\Lambda_S$ is significantly different from zero.

Crucially, this construction imposes \emph{no assumption} that all settings share a single data-generating process.
The settings can be an arbitrarily eclectic mixture---public goods games, dictator games, or entirely unrelated tasks.   
Inference remains valid even when the model's performance differs sharply across sub-domains; the variability of estimates widens (or narrows) in proportion to the observed heterogeneity.

The remainder of this section is devoted to implementing the above framework on a pre-committed family of strategic games.
This set comprises \TotalGames{} novel and unique permutations of \citeauthor{1120Arad2012}'s money request game.
We randomly sample \NGamesS{} of these games for human subjects and \aisub{s} to play in a preregistered experiment.
We use this data to estimate the relative capacity of different \aisub{s} to predict human responses at scale.
In particular, we return to the strategic sample of optimized level-$k$ \aisub{s} presented in Table~\ref{tab:persona} from Section~\ref{sec:predict_new_AR}.
We compare these agents' ability to predict human responses across the \NGamesS{} games to the baseline AI, a cognitive hierarchy model, and symmetric Nash equilibria.
Because the games (and human subjects) are randomly sampled from the population according to a known distribution, confidence intervals over the comparisons are valid over the \TotalGames{} games.
